% !TeX root = ../sections/03_solutions.tex

\subsection{1_1}
\begin{solution}
	周期 \(2\pi\) で拡張された関数として考える。
	
	まず，
	\[
	\frac{1}{x}
	\]
	は \(x=0\) で定義されず，さらに
	\[
	\lim_{x\to 0+}\frac{1}{x}=+\infty,
	\qquad
	\lim_{x\to 0-}\frac{1}{x}=-\infty
	\]
	である。したがって，有限の片側極限をもつ不連続ではないので，
	区分的に連続ではない。よって，滑らかでも区分的に滑らかでもない。
	
	次に，
	\[
	\sqrt{|x|}
	\]
	は \(x=0\) でも連続であり，端点においても
	\[
	\sqrt{|-\pi|}=\sqrt{\pi}
	\]
	であるから，周期拡張しても連続である。
	しかし \(x=0\) で導関数が発散するため，滑らかではない。
	また，通常の意味での区分的に滑らかな関数，すなわち導関数まで区分的に連続である関数とはみなさない。
	
	次に，
	\[
	x
	\]
	は区間 \((-\pi,\pi)\) では滑らかであるが，周期拡張すると
	\[
	\lim_{x\to \pi-0}x=\pi,
	\qquad
	f(-\pi)=-\pi
	\]
	となるので，端点でジャンプ不連続をもつ。
	したがって連続ではないが，区分的に連続であり，区分的に滑らかである。
	
	次に，
	\[
	x^2
	\]
	は区間内部で滑らかであり，端点でも
	\[
	(-\pi)^2=\pi^2
	\]
	となるため，周期拡張しても連続である。
	しかし導関数は
	\[
	\frac{d}{dx}x^2=2x
	\]
	であり，
	\[
	\lim_{x\to \pi-0}2x=2\pi,
	\qquad
	\lim_{x\to -\pi+0}2x=-2\pi
	\]
	となるので，周期関数としては滑らかではない。
	ただし，有限個の点を除けば滑らかなので，区分的に滑らかである。
	
	最後に，
	\[
	\sin 3x
	\]
	は三角関数であり，周期拡張しても値や導関数が自然につながる。
	したがって滑らかである。
	
	以上より，
	\[
	\begin{array}{c|l}
		\text{分類} & \text{該当する関数} \\ \hline
		\text{連続な関数}
		&
		\sqrt{|x|},\ x^2,\ \sin 3x
		\\[2mm]
		\text{区分的に連続な関数}
		&
		\sqrt{|x|},\ x,\ x^2,\ \sin 3x
		\\[2mm]
		\text{滑らかな関数}
		&
		\sin 3x
		\\[2mm]
		\text{区分的に滑らかな関数}
		&
		x,\ x^2,\ \sin 3x
	\end{array}
	\]
	である。
\end{solution}